\setcounter{equation}{0}
\section{Library Documentation}

This section documents the simulator API that would replace the historical 2004 library if the sparse-gate strategy of Section~\ref{sec:parallel} were implemented in full. The 2004 source bundled with this distribution (under \texttt{implementation/\allowbreak original/}) implements the older dense-matrix approach and is preserved as a historical artifact; the API below is what the corrected design calls for.

\subsection{The quantum register}

A quantum register is represented by an opaque struct holding a local slice of the distributed state vector together with the MPI bookkeeping that identifies which slice this process owns:

\begin{lstlisting}
typedef struct {
    complex double *amp;       /* local slice of the state vector  */
    int    n_qubits;           /* total number of qubits (global)  */
    int    rank, n_procs;      /* MPI rank and process count       */
    size_t local_size;         /* 2^n_qubits / n_procs amplitudes  */
} qreg;
\end{lstlisting}

\paragraph{Construction.} \verb|qreg_create| allocates a register on the given MPI communicator and zero-initialises the amplitude array:
\begin{lstlisting}
qreg *q = qreg_create(20, MPI_COMM_WORLD);
\end{lstlisting}

\paragraph{Initialisation.} \verb|qreg_init_basis| sets the register to a definite computational basis state $\ket{k}$:
\begin{lstlisting}
qreg_init_basis(q, 0);   /* sets the register to |0...0> */
qreg_init_basis(q, 3);   /* sets it to |...011>          */
\end{lstlisting}

\paragraph{Destruction.} \verb|qreg_destroy| frees all owned memory.

\subsection{Gate primitives}

Single-qubit gates apply a $2\times 2$ unitary in place on the local amplitude slice, using the sparse-iteration pattern of Section~\ref{sec:parallel}. The library provides direct wrappers for the standard one-qubit gates and a generic primitive for user-supplied $2\times 2$ matrices:

\begin{lstlisting}
void apply_h     (qreg *q, int target);                 /* Hadamard */
void apply_x     (qreg *q, int target);                 /* Pauli-X  */
void apply_y     (qreg *q, int target);                 /* Pauli-Y  */
void apply_z     (qreg *q, int target);                 /* Pauli-Z  */
void apply_phase (qreg *q, int target, double theta);   /* diag(1,e^{i theta}) */
void apply_u     (qreg *q, int target, complex double u[2][2]);
\end{lstlisting}

Two-qubit controlled gates take a control qubit and a target qubit:
\begin{lstlisting}
void apply_cnot            (qreg *q, int control, int target);
void apply_controlled_phase(qreg *q, int control, int target, double theta);
void apply_cu              (qreg *q, int control, int target,
                            complex double u[2][2]);
void apply_swap            (qreg *q, int a, int b);
\end{lstlisting}

For each call, the library inspects whether the target/control qubits are local to this process. If both are local (the common case --- see Section~\ref{sec:parallel}) the gate is applied with no communication; otherwise the library transparently performs the pairwise amplitude exchange and then the local update.

\subsection{Measurement and inspection}

\verb|measure_all| collapses the state and returns the observed basis state on every process:
\begin{lstlisting}
size_t outcome = measure_all(q);
\end{lstlisting}

\verb|prob_of| returns $|\langle k | \psi \rangle|^2$ for a chosen basis state $k$ without disturbing the register:
\begin{lstlisting}
double p = prob_of(q, 42);
\end{lstlisting}

\verb|measure_qubit| performs a partial measurement of a single qubit and projects the state accordingly:
\begin{lstlisting}
int bit = measure_qubit(q, 3);
\end{lstlisting}

\subsection{Higher-level routines}

The chapters on the QFT (Section~\ref{sec:qft}), Grover (Section~\ref{sec:grover}), and Shor (Section~\ref{sec:shor}) build on the primitives above. The QFT is a double loop of Hadamards and controlled phases; Grover is a uniform-superposition initialisation followed by alternating oracle and diffusion calls; Shor combines phase estimation, controlled modular exponentiation, and inverse QFT. Each is exposed as a top-level function in the library so applications need not assemble the circuit by hand:

\begin{lstlisting}
void apply_qft         (qreg *q, int start, int n);
void apply_qft_inverse (qreg *q, int start, int n);
void apply_grover      (qreg *q, int n, oracle_fn oracle, int iterations);
void apply_shor_period (qreg *q, int counting_start, int t,
                        int target_start, int n,
                        int a, int N);
\end{lstlisting}

\subsection{Note on the bundled 2004 implementation}

\begin{sloppypar}
The C source under \texttt{implementation/original/} (files \texttt{matrix.c}, \texttt{parallel.c}, \texttt{qubit.c}, \texttt{standart.c}) implements a different API based on a dense \verb|matrix| struct of type \verb|__complex__ double**|:
\end{sloppypar}

\begin{lstlisting}
typedef struct {
    int row_size;
    int column_size;
    __complex__ double ** value;
} matrix;
\end{lstlisting}

with operations \verb|create_matrix|, \verb|tensor_product|, \verb|dot_product|, and the MPI primitives \verb|send_matrix|, \verb|get_matrix|, and \verb|broadcast_matrix|. This code targets the original dense-matrix strategy critiqued in Section~\ref{sec:parallel} and is retained here for historical reference. The API documented above is the design that should replace it; implementation of the corrected library is left as future work.
